\documentclass[12pt,a4paper]{article}
\usepackage[margin=2.5cm]{geometry}
\usepackage{graphicx}
\usepackage{booktabs}
\usepackage{enumitem}
\usepackage{hyperref}
\usepackage{xcolor}
\usepackage{titlesec}
\usepackage{fancyhdr}
\usepackage{tabularx}

\pagestyle{fancy}
\fancyhf{}
\rhead{Implementation Progress Report}
\lhead{Solana Trading Bot}
\rfoot{Page \thepage}

\title{\textbf{Solana Autonomous Trading Bot}\\[0.3cm]\large Implementation Progress Report}
\author{Trading Bot Project}
\date{April 2026}

\begin{document}

\maketitle
\thispagestyle{empty}

\section{Introduction}

This report presents the current implementation status of our autonomous Solana trading bot, a 14-phase project designed to detect and execute on-chain trading opportunities in real time. The system is built on a hybrid architecture combining a TypeScript backend for on-chain execution, a Python AI server for predictive signals, Anchor smart contracts deployed on Solana DevNet, and a React frontend dashboard for multi-user management.

The pipeline, illustrated in Figure~\ref{fig:pipeline}, follows a three-stage flow: \textbf{Opportunity Discovery} (Layer~1), \textbf{Execution} (Layer~2), and \textbf{Protection} (Layer~3). At its core, the \texttt{TradingEngine.tick()} method orchestrates a single pass through the entire pipeline every 10 seconds.

\begin{figure}[h]
\centering
\includegraphics[width=\textwidth]{pipeline.png}
\caption{Trading bot pipeline architecture showing the fast path (flash loans), slow path (own capital + AI), verification stages, and the planned three-layer agentic pipeline.}
\label{fig:pipeline}
\end{figure}

\section{Pipeline Architecture}

\subsection{Opportunity Sources}

The system discovers opportunities from two categories of sources:

\begin{itemize}[nosep]
    \item \textbf{Pure Code detectors} (fully implemented): Arbitrage detection via triangular constant-product AMM formula, liquidation hunting for undercollateralized positions, and yield rate monitoring across pools.
    \item \textbf{AI Model signals} (implemented): Sentiment analysis (Gemini 2.0 Flash with VADER fallback), GRU regime prediction (crash/ranging/trending classification), XGBoost direction forecasting, and whale movement tracking.
\end{itemize}

\subsection{Path Classifier}

The \texttt{OpportunityRouter} classifies each discovered opportunity into one of two execution paths:

\begin{itemize}[nosep]
    \item \textbf{Fast Path (Flash Loan)}: For arbitrage and liquidation opportunities requiring zero own capital. The transaction atomically borrows, executes swaps, and repays within a single Solana transaction. This path is \textbf{fully operational} on DevNet.
    \item \textbf{Slow Path (Own Capital)}: For yield farming and AI-driven opportunities that require the user's deposited capital. This path uses the \texttt{DecisionModel} fusion layer, which combines regime prediction, volatility estimation, sentiment scoring, and whale signals into a unified confidence score before committing capital.
\end{itemize}

\subsection{Verification Pipeline}

Before any opportunity reaches execution, it must pass through three verification stages:

\begin{enumerate}[nosep]
    \item \textbf{Threshold check}: Minimum profit threshold after fees.
    \item \textbf{S1--S4 Safety Pipeline}: A sequential four-stage filter:
    \begin{itemize}[nosep]
        \item S1: On-chain heuristics (mint authority revoked, dev wallet percentage, token age)
        \item S2: Honeypot detection (simulated sell via pool existence)
        \item S3: Isolation forest anomaly scoring (Redis-cached, threshold $< 0.7$)
        \item S4: Anomaly detection (holder count, wash trading patterns)
    \end{itemize}
    \item \textbf{Protection Manager}: Auto-pause status, trading hours enforcement, daily drawdown limits, and slippage guards.
\end{enumerate}

All four safety stages and the protection manager are \textbf{fully implemented and tested} on DevNet.

\subsection{Execution Layer}

Verified opportunities are executed via Anchor smart contracts deployed on Solana DevNet:

\begin{itemize}[nosep]
    \item \textbf{AMM Program}: Supports \texttt{swap}, \texttt{add\_liquidity}, and \texttt{remove\_liquidity} instructions.
    \item \textbf{Flash Loan Program}: Implements \texttt{flash\_borrow} and \texttt{flash\_repay} with instruction introspection to enforce atomicity (repayment must exist in the same transaction).
    \item \textbf{Per-User Vault Program}: Anchor program written but \textbf{not yet deployed}. Will enable individual user capital management.
\end{itemize}

\subsection{Planned: Three-Layer Agentic Pipeline}

As shown in the lower portion of the pipeline diagram, the system is designed to eventually incorporate a three-layer agentic architecture:

\begin{itemize}[nosep]
    \item \textbf{Analyzer Agent}: Market analysis and signal interpretation.
    \item \textbf{Strategist Agent}: Strategy selection and parameter optimization.
    \item \textbf{Executor Agent}: Trade execution with customization per user preferences.
\end{itemize}

This agentic layer is \textbf{not yet started} and represents a future enhancement beyond the current 14-phase plan.

\section{Current Implementation Status}

\subsection{What Works End-to-End Today}

The fast path pipeline is fully operational on DevNet:

\begin{enumerate}[nosep]
    \item The engine reads on-chain pool reserves and detects triangular arbitrage opportunities.
    \item Each opportunity passes through the S1--S4 safety pipeline.
    \item The protection manager validates trading conditions (hours, drawdown, pause state).
    \item A flash loan transaction is built atomically: \texttt{flash\_borrow} $\rightarrow$ 3 AMM swaps $\rightarrow$ \texttt{flash\_repay}.
    \item The transaction is submitted and profit is recorded.
\end{enumerate}

This has been validated by the end-to-end test script (\texttt{test\_e2e\_pipeline.ts}), which boots the engine, runs a single tick, and asserts seven correctness checks.

\subsection{Multi-User Platform (Phase 13)}

The multi-user architecture is operational with the following components:

\begin{itemize}[nosep]
    \item \textbf{Backend API}: Express server on port 3001 with REST endpoints for user management, engine control (start/stop), status polling, and configuration updates. WebSocket support for real-time tick event streaming.
    \item \textbf{User Registry}: Each user gets an independent \texttt{TradingEngine} instance. Configuration changes propagate live to the engine's protection manager and opportunity filters.
    \item \textbf{Frontend Dashboard}: React + Vite application with Solana wallet adapter integration (Phantom, Solflare). Pages include a live cockpit with KPI display, a bot builder wizard for configuration, and analytics views.
    \item \textbf{Pending}: Per-user vault smart contract deployment (program written, not deployed), and the vault deposit/withdraw UI exists in the frontend but is non-functional until the program is live.
\end{itemize}

\section{AI Model Integration}

The AI server is a FastAPI application exposing prediction endpoints consumed by the TypeScript backend. The following models are integrated:

\begin{itemize}[nosep]
    \item \textbf{GRU Regime Predictor}: A recurrent neural network trained on 18 market features with a sequence length of 24, classifying the current market regime into three classes: \emph{crash}, \emph{ranging}, or \emph{trending}. The regime prediction feeds into the \texttt{DecisionModel} to gate slow-path trades --- the bot avoids committing own capital during crash regimes.
    \item \textbf{XGBoost Direction Classifier}: A gradient-boosted tree model predicting short-term price direction. Combined with regime context, it provides a directional confidence score used in the decision fusion layer.
    \item \textbf{Sentiment Analysis}: Uses the Gemini 2.0 Flash API for LLM-based market sentiment extraction from news and social data. Falls back to VADER (a lexicon-based sentiment analyzer) when the API key is unavailable or rate-limited. Scores are cached in Redis with configurable TTL.
    \item \textbf{Whale Movement Tracker}: Monitors large wallet balance changes on-chain and converts them into directional signals (accumulation vs.\ distribution), tracked per user configuration.
\end{itemize}

The \texttt{DecisionModel} fusion layer aggregates these signals --- regime, direction, sentiment, and whale activity --- into a single confidence score. Slow-path opportunities are only executed when this score exceeds the user-configured threshold, ensuring that own-capital trades are backed by multi-model consensus.

\section{Conclusion}

The trading bot pipeline is fully operational on Solana DevNet, covering both execution paths. The fast path executes atomic flash loan arbitrage with zero own capital, while the slow path leverages a multi-model AI fusion layer (GRU regime prediction, XGBoost direction, LSTM price signals, sentiment analysis, and whale tracking) to make informed capital allocation decisions. The four-stage safety pipeline, protection manager, and Anchor smart contracts (AMM + flash loan) are all deployed and tested end-to-end.

The multi-user platform extends this with an Express + WebSocket API server and a React frontend dashboard, giving each user an independent engine instance with live configuration propagation. The planned three-layer agentic pipeline (Analyzer, Strategist, Executor) represents the next architectural evolution, adding autonomous decision-making capabilities on top of the existing infrastructure. Remaining work focuses on per-user vault deployment for capital isolation and mainnet hardening.

\end{document}
